\documentclass[conference]{IEEEtran}

\usepackage{amsmath}
\usepackage{booktabs}
\usepackage{cite}
\usepackage{url}

\begin{document}

\title{When FLOPs Mislead: Benchmarking Neural Network Inference Energy on Apple Silicon}

\author{
\IEEEauthorblockN{Aryan Shah}
\IEEEauthorblockA{Texas Academy of Mathematics and Science\\
University of North Texas\\
Denton, TX, USA\\
aryanshah.2431@gmail.com}
}

\maketitle

\begin{abstract}
Energy is a deployment constraint for edge-AI systems, but neural-network
efficiency is still often discussed using architecture-level proxies such as
FLOPs, parameter count, and model size. We first reproduce a five-model
Apple-Silicon CPU-only PyTorch comparison, then add a new direct
\texttt{powermetrics} audit over 17 architecture variants on an Apple M4 Pro.
The new audit contains 170 measured windows, 30 seconds per window, raw power
logs, per-window inference counts, confidence intervals, and full environment
metadata. Three follow-up audits then test the result under MPS/Metal
execution, CPU batch size 4, and a smaller $128\times128$ CPU input. Across
the 17 directly measured CPU batch-1 architectures, latency explains
nearly all between-architecture energy variation ($R^2=0.9925$), while FLOPs
is much weaker ($R^2=0.4008$) and parameter count is weaker still
($R^2=0.0977$). The expanded data preserve the central five-model result:
runtime behavior, not nominal arithmetic count, determines energy ordering
under this fixed CPU backend, and the sensitivity checks show which parts of
the ranking persist when the backend, batch size, or input size changes. For
example, ConvNeXt-micro has only 0.448
GFLOPs but consumes slightly more energy than a 3.414-GFLOP residual CNN
because both run at roughly the same latency. A residualized conditional
WGAN-GP is retained only as supplemental sparse-trial modeling. The results
support a practical benchmarking rule for edge inference: report direct
latency and auditable power traces under the deployment runtime instead of
ranking models by FLOP count alone.
\end{abstract}

\begin{IEEEkeywords}
edge AI, benchmarking, neural network inference, energy measurement, Apple
Silicon, energy-delay product
\end{IEEEkeywords}

\section{Introduction}

High-performance embedded and edge systems increasingly run neural-network
inference under strict energy and latency budgets. Model selection is often
guided by FLOPs, parameter count, or model size because these quantities are
easy to compute before deployment. They are useful descriptors, but they do
not directly measure wall-clock work, memory traffic, kernel efficiency, or
the behavior of a specific runtime on a specific processor \cite{horowitz2014,
eyeriss2016, mobilenetv3}.

This paper studies a narrow but practical question: for a fixed Apple Silicon
CPU-only PyTorch runtime, which common architecture proxy best tracks measured
inference energy? The answer in this benchmark is latency, not FLOPs. This is
important for HPEC-style edge deployments because the relevant cost is the
energy actually paid by the hardware under the deployment runtime, not the
nominal arithmetic count of the neural network.

The contributions are:
\begin{itemize}
    \item A five-architecture Apple-Silicon inference-energy benchmark
    comparing MobileNetV3-Small, MobileNetV2, ResNet-18, Tiny-ViT-5M, and
    EfficientNet-B0.
    \item A new 17-architecture Apple M4 Pro direct-energy audit with 170
    \texttt{powermetrics} windows, raw logs, per-window inference counts,
    confidence intervals, and environment metadata.
    \item Three targeted follow-up audits for limitations: MPS/Metal backend
    comparison, CPU batch-size sensitivity at batch 4, and CPU input-size
    sensitivity at $128\times128$.
    \item A compact analysis showing that latency is a much stronger predictor
    of energy than FLOPs in both the original five-model table and the
    expanded 17-architecture direct-energy table.
    \item A scoped synthetic-trial modeling note: residualized conditional
    GANs can support exploratory repeated-trial analysis, but they should not
    replace independent power measurements.
\end{itemize}

\section{Related Work}

Efficient neural-network design has produced architectures such as MobileNet,
MobileNetV2, MobileNetV3, EfficientNet, SqueezeNet, and ShuffleNet, which
reduce arithmetic through depthwise separable convolutions, inverted
residuals, compound scaling, and channel grouping \cite{mobilenet,
mobilenetv2, mobilenetv3, efficientnet, squeezenet, shufflenet}. These
designs often improve accuracy/FLOP tradeoffs, but hardware energy depends on
more than FLOPs.

Prior systems work shows that data movement and memory hierarchy behavior can
dominate arithmetic energy \cite{horowitz2014, eyeriss2016}. Hardware-aware
model design similarly emphasizes that latency and hardware utilization can
diverge from FLOP counts \cite{mobilenetv3, efficientnet}. Green AI and ML
systems benchmarking motivate reporting measured deployment costs rather than
abstract complexity alone \cite{greenai, carbontracker}. This paper adds a
small, auditable Apple-Silicon case study centered on per-inference energy.

\section{Methodology}

\subsection{Runtime and Hardware Scope}

All primary results use CPU-only PyTorch inference with batch size 1 and input
shape $1\times3\times224\times224$. CPU-only execution was chosen to isolate a
portable runtime path and avoid Neural Engine or Metal scheduling effects. It
does not claim to be the most energy-efficient Apple Silicon backend.

The repository contains multiple Apple-Silicon energy sources. The first is the
paper-aligned original five-model mean table, which reports one mean latency
and one mean energy per model from Apple-Silicon CPU inference measured with
\texttt{powermetrics} at 1 Hz. The raw logs for that original table are
unavailable, and the local artifact does not identify it as M1, so we do not
infer confidence intervals or a specific chip generation for it.

The main new evidence is a 17-architecture Apple M4 Pro audit collected in
\texttt{measured\_energy\_powermetrics\_m4pro\_arch17\_cpu}. It uses 10
30-second windows per architecture, 170 measured windows total, 1 Hz
\texttt{powermetrics} sampling, 20 warmup inferences before each window,
randomized architecture order, one CPU thread, batch size 1, and input shape
$1\times3\times224\times224$. The audit preserves raw \texttt{powermetrics}
text logs, per-window inference counts, per-window energy estimates, summary
statistics, confidence intervals, and a JSON environment record. Earlier
five-model M4 and retained M5 summaries are treated as supporting audits, not
merged with the new 17-architecture table.

After the main audit, three shorter follow-up audits were collected on the
same M4 Pro environment to address backend and input-scope limitations. The
first repeats the 17-architecture suite with PyTorch MPS/Metal at batch size 1
and $224\times224$ input. The second repeats the CPU suite at batch size 4 and
$224\times224$ input. The third repeats the CPU suite at $128\times128$ input
for 15 architectures; the two patch-mixer variants are excluded from this
input-size check because their token MLPs are shape-fixed to the $224\times224$
registry definition. These follow-up audits use five 20-second windows per
architecture and are reported separately from the main CPU batch-1 table.

\begin{table*}[t]
\centering
\caption{Direct repeated Apple-Silicon energy-audit environments.}
\label{tab:environment}
\begin{tabular}{lll}
\toprule
Component & New 17-architecture M4 Pro audit & Retained M5 Pro audit \\
\midrule
Device & MacBook Pro Mac16,8 & MacBook Pro Mac17,9 \\
Chip & Apple M4 Pro, 12 cores & Apple M5 Pro, 15 cores \\
Memory & 24 GB RAM & 24 GB RAM \\
OS & macOS 26.2 build 25C56 & macOS 26.4.1 \\
Python/PyTorch & 3.9.6 / 2.8.0 & 3.12.1 / 2.3.1 \\
torchvision/timm & 0.23.0 / 1.0.24 & 0.18.1 / 1.0.24 \\
Threads/input & 1 thread, batch 1 & 1 thread, batch 1 \\
Power logging & \texttt{powermetrics}, 1 Hz & \texttt{powermetrics}, 1 Hz \\
Protocol & 10 windows/model, 30 s/window & 10 windows/model, 20 s/window \\
\bottomrule
\end{tabular}
\end{table*}

\subsection{Metrics}

We report per-inference energy $E$ in joules, latency $L$ in milliseconds, and
average power $P=1000E/L$. Because power is sampled once per second, energy is
interpreted as window-level energy divided by the number of inferences
completed in the window. We also report Energy-Delay Product (EDP),
\begin{equation}
    \mathrm{EDP} = E \cdot L.
\end{equation}
EDP is not accuracy-normalized; it measures energy-latency cost for fixed
input-shape inference.

\subsection{Benchmark Procedure}

Each model is instantiated once, switched to evaluation mode, and run with
gradients disabled. Warmup inferences are discarded so that module loading and
one-time setup do not dominate short windows. During each measured window, the
script repeatedly executes forward passes until the window expires, records
the number of completed inferences, and divides the package-energy estimate
for that same window by the inference count. Raw \texttt{powermetrics} output
is retained, so the summary CSV is not the only source of evidence.

This windowed procedure is deliberately conservative for sub-100 ms inference.
It does not try to attribute a single instantaneous power sample to one
forward pass. Instead, it measures many repeated inferences under a controlled
loop and reports the average energy paid by that loop.

\subsection{Supplemental Synthetic Modeling}

The repository includes a conditional WGAN-GP pipeline used to explore
synthetic repeated trials under sparse measurement budgets. An initial
multi-device raw-scale model mode-collapsed. The revised model uses
within-combination residuals:
\begin{equation}
    \tilde{x} = \mathrm{clip}\left((x-\mu_c)/\sigma_c,-3,3\right).
\end{equation}
Generated residuals are mapped back using the same $\mu_c$ and $\sigma_c$.
This synthetic component is treated as supplemental. It can reproduce
grounded means and rankings, but its spreads are not a substitute for
held-out measured trial-level energy data.

\section{Results}

Table~\ref{tab:original} summarizes the five-model paper-aligned original mean
table. The key result is that energy follows latency rather than FLOPs.
EfficientNet-B0 has only 0.39 GFLOPs, but it has the largest latency and
energy. ResNet-18 has 1.82 GFLOPs, yet consumes less than half the energy of
EfficientNet-B0 in this setup.

\begin{table}[t]
\centering
\caption{Paper-aligned original five-model mean table. EDP is computed from
measured energy and latency.}
\label{tab:original}
\begin{tabular}{lrrrr}
\toprule
Model & FLOPs & Lat. & Energy & EDP \\
 & (G) & (ms) & (J) & (J ms) \\
\midrule
MobileNetV3-S & 0.22 & 15.14 & 0.080 & 1.211 \\
MobileNetV2 & 0.30 & 20.07 & 0.106 & 2.127 \\
ResNet-18 & 1.82 & 20.68 & 0.110 & 2.275 \\
Tiny-ViT-5M & 1.30 & 41.33 & 0.219 & 9.051 \\
EfficientNet-B0 & 0.39 & 55.08 & 0.292 & 16.083 \\
\bottomrule
\end{tabular}
\end{table}

\subsection{Predictor Analysis}

We fit univariate linear models predicting energy from latency, FLOPs, and
parameter count in the original mean table. Latency explains nearly all
between-model energy variation ($R^2=0.99999$), while FLOPs explains almost
none ($R^2=0.00008$). Leave-one-out checks preserve the same qualitative
result: latency $R^2$ remains above 0.99997, while FLOPs remains weak.

\subsection{Seventeen-Architecture M4 Pro Energy Audit}

Table~\ref{tab:m4arch17} reports the new direct M4 Pro audit. Each architecture
has 10 measured 30-second windows, with raw logs and per-window rows stored in
the artifact directory. Confidence intervals are Student-$t$ 95\% intervals
over the 10 measured windows. This table replaces the earlier latency-only
17-architecture sweep as the main expanded result: the same architecture
families now have direct \texttt{powermetrics} energy measurements rather than
constant-power energy proxies.

\begin{table*}[t]
\centering
\caption{New direct \texttt{powermetrics} audit over 17 architecture variants on Apple M4 Pro. Rows are sorted by measured energy.}
\label{tab:m4arch17}
\small
\begin{tabular}{lrrrr}
\toprule
Architecture & FLOPs & Lat. & Energy & 95\% CI \\
 & (G) & (ms) & (J) & (J) \\
\midrule
global\_avg\_mlp\_tiny & 0.000 & 0.021 & 0.00012 & [0.0001,0.0001] \\
patch\_mixer\_small & 0.042 & 0.431 & 0.00284 & [0.0025,0.0032] \\
cnn\_tiny\_2conv & 0.040 & 0.770 & 0.00502 & [0.0045,0.0055] \\
patch\_mixer\_medium & 0.101 & 0.741 & 0.00534 & [0.0044,0.0063] \\
depthwise\_small & 0.023 & 1.289 & 0.00797 & [0.0074,0.0086] \\
cnn\_small\_4conv & 0.170 & 1.306 & 0.00955 & [0.0078,0.0113] \\
grouped\_conv\_cnn & 0.171 & 1.746 & 0.01120 & [0.0103,0.0121] \\
depthwise\_medium & 0.050 & 1.904 & 0.01328 & [0.0113,0.0153] \\
cnn\_medium\_6conv & 0.746 & 2.414 & 0.01587 & [0.0129,0.0189] \\
attention\_gate\_cnn & 0.526 & 2.390 & 0.01663 & [0.0145,0.0187] \\
squeeze\_expand\_cnn & 0.597 & 3.089 & 0.01955 & [0.0185,0.0206] \\
residual\_cnn\_small & 1.062 & 3.434 & 0.02231 & [0.0197,0.0249] \\
bottleneck\_residual & 0.724 & 3.459 & 0.02252 & [0.0203,0.0248] \\
inverted\_residual\_small & 0.279 & 3.311 & 0.02313 & [0.0202,0.0261] \\
residual\_cnn\_medium & 3.414 & 6.757 & 0.04259 & [0.0387,0.0464] \\
convnext\_micro & 0.448 & 6.832 & 0.04417 & [0.0418,0.0466] \\
inverted\_residual\_wide & 0.812 & 6.811 & 0.04898 & [0.0392,0.0588] \\
\bottomrule
\end{tabular}
\end{table*}

The expanded direct audit supports the same qualitative conclusion as the
original five-model table. Across 17 architectures, latency explains nearly
all measured energy variation ($R^2=0.9925$), while FLOPs is much weaker
($R^2=0.4008$) and parameter count is weaker still ($R^2=0.0977$). The
highest-energy group is not simply the highest-FLOP group: ConvNeXt-micro uses
0.448 GFLOPs but consumes slightly more energy than the 3.414-GFLOP residual
CNN medium because both run near 6.8 ms per inference. Inverted-residual-wide
is the highest-energy architecture despite using less than one GFLOP. These
counterexamples show why FLOPs alone is not a reliable deployment-energy proxy
for this CPU backend.

\begin{table}[t]
\centering
\caption{Predictors of measured energy across the main evidence sources.}
\label{tab:predictors}
\begin{tabular}{lrr}
\toprule
Dataset and predictor & $R^2$ & Pearson $r$ \\
\midrule
Original: latency & 0.99999 & 0.99999 \\
Original: FLOPs & 0.00008 & -0.009 \\
April M4 audit: latency & 0.980 & 0.990 \\
April M4 audit: FLOPs & 0.001 & 0.034 \\
M4 17-architecture: latency & 0.9925 & 0.9963 \\
M4 17-architecture: FLOPs & 0.4008 & 0.6331 \\
M4 17-architecture: params & 0.0977 & 0.3125 \\
M4 MPS batch 1: latency & 0.6236 & 0.7897 \\
M4 MPS batch 1: FLOPs & 0.7892 & 0.8884 \\
M4 CPU batch 4: latency & 0.9914 & 0.9957 \\
M4 CPU batch 4: FLOPs & 0.4397 & 0.6631 \\
M4 CPU $128^2$: latency & 0.9986 & 0.9993 \\
M4 CPU $128^2$: FLOPs & 0.1784 & 0.4223 \\
Retained M5 audit: latency & 0.975 & -- \\
Retained M5 audit: FLOPs & 0.005 & -- \\
\bottomrule
\end{tabular}
\end{table}

\subsection{Backend, Batch, and Input-Size Checks}

Table~\ref{tab:sensitivity} summarizes the three targeted follow-up audits.
The CPU batch-size and input-size checks preserve the central predictor
result: latency remains above $R^2=0.99$, while FLOPs remains much weaker. The
CPU batch-4 ranking is also very close to the main CPU batch-1 ranking
(Spearman $\rho=0.990$ over 17 shared architectures), and the $128\times128$
CPU ranking remains close over the 15 shared architectures
(Spearman $\rho=0.993$).

MPS/Metal changes the picture. It is much faster and lower-energy for most
larger convolutional and residual architectures, but not for the smallest
models where dispatch overhead dominates. The MPS energy ranking is still
related to the CPU ranking (Spearman $\rho=0.914$), but FLOPs becomes the
stronger single predictor in this short backend audit ($R^2=0.7892$ for FLOPs
versus $R^2=0.6236$ for latency). This reinforces the main methodological
point: direct backend-specific measurement is required, because one runtime's
proxy relationship should not be assumed to transfer unchanged to another
runtime.

\begin{table*}[t]
\centering
\caption{Targeted M4 Pro follow-up audits. No measured windows are filtered or winsorized.}
\label{tab:sensitivity}
\begin{tabular}{lrrrrrr}
\toprule
Audit & Models & Windows & Lat. $R^2$ & FLOPs $R^2$ & Energy-rank $\rho$ & Largest retained window \\
\midrule
CPU batch 1, $224^2$ & 17 & 170 & 0.9925 & 0.4008 & 1.000 & inverted\_residual\_wide, 0.0865 J \\
MPS batch 1, $224^2$ & 17 & 85 & 0.6236 & 0.7892 & 0.914 & residual\_cnn\_medium, 0.0273 J \\
CPU batch 4, $224^2$ & 17 & 85 & 0.9914 & 0.4397 & 0.990 & residual\_cnn\_medium, 0.2685 J \\
CPU batch 1, $128^2$ & 15 & 75 & 0.9986 & 0.1784 & 0.993 & inverted\_residual\_wide, 0.0243 J \\
\bottomrule
\end{tabular}
\end{table*}

\subsection{Interpretation for Device Expansion}

The original mean table, the earlier five-model M4 audit, the new
17-architecture M4 audit, the three M4 follow-up audits, and the retained M5
summary are intentionally kept
separate. The original table supports the earlier paper-aligned comparison;
the new M4 audit supplies the strongest expanded direct-energy evidence with
raw logs and repeated-window uncertainty; the follow-up audits test whether
the main pattern survives backend, batch-size, and input-size changes; and the
M5 result is retained as a cross-device audit summary from the prior HPEC
artifact. Agreement between
them is therefore evidence for robustness of the qualitative pattern, not
identity of absolute energy values.

This distinction matters for future device expansion. Apple Silicon CPU-only
\texttt{powermetrics} data, Windows CPU data, and Windows GPU data should be
reported as separate backend/device conditions. In particular, Windows GPU
power logs from tools such as \texttt{nvidia-smi} are not directly equivalent
to Apple package-power estimates. They can strengthen the paper if collected
with matched model inputs, synchronization around GPU timing, raw power logs,
and environment metadata, but they should not be merged into a single
undifferentiated energy claim.

\section{Discussion}

The empirical pattern is consistent across the available Apple-Silicon energy
sources: latency is a far stronger proxy for energy than FLOPs. This
does not mean latency is a physical cause of energy by itself. Rather, under a
fixed runtime and input shape, latency captures many implementation details
that FLOPs misses: memory traffic, operator fusion, dispatch overhead,
vectorization, cache behavior, and kernel availability.

The largest counterexamples to FLOP-based reasoning now appear in the expanded
M4 audit. ConvNeXt-micro has 7.6$\times$ fewer FLOPs than residual-CNN-medium
but slightly higher measured energy because their latencies are almost the
same. Inverted-residual-wide has 0.812 GFLOPs, far below the 3.414 GFLOPs of
residual-CNN-medium, but it is the most energy-consuming architecture in the
17-model table. FLOPs alone would therefore make the wrong deployment
recommendation for this fixed CPU runtime.

The M4 audits are also valuable because they add uncertainty information absent
from the original table. High-power windows are retained rather than removed,
and each table entry links back to raw text logs. This makes the evidence more
auditable and prevents synthetic data from being confused with physical
measurement.

\section{Limitations}

The new audit improves the main scope limitation by collecting direct energy
for all 17 architecture variants and adding targeted MPS/Metal, batch-size,
and input-size checks. However, the main result remains a local Apple M4 Pro
study, and the follow-up audits are shorter than the primary 170-window CPU
batch-1 audit. The $128\times128$ input-size check also covers 15
architectures rather than all 17 because the patch-mixer registry entries are
shape-fixed.

The paper still does not evaluate Apple Neural Engine, quantized deployment
paths, CUDA, or other vendor-specific GPU inference backends. Those conditions
require separate timing synchronization and power logging protocols. The
\texttt{powermetrics} sampling interval is also coarse relative to
per-inference latency, so energy is attributed at the window level rather than
at the level of one forward pass. The local rebuild does not include a raw M5
Pro artifact directory, so the M5 entry is reported only as a retained prior
HPEC audit summary. Finally, the residualized WGAN-GP is evaluated against
grounded training
means and should be interpreted as a supplemental modeling tool, not as proof
of physical repeated-trial variance.

\section{Conclusion}

On Apple Silicon under CPU-only PyTorch inference, per-inference energy cannot
be reduced to FLOPs or parameter count. In the original five-model mean table,
latency tracks energy at $R^2=0.99999$, while FLOPs explains almost none of
the measured variation. The new 17-architecture Apple M4 Pro
\texttt{powermetrics} audit strengthens that result with broader direct energy
coverage: latency reaches $R^2=0.9925$, while FLOPs reaches only
$R^2=0.4008$. For embedded and edge inference, this supports a practical
measurement standard: report latency and auditable power traces under the
deployment runtime instead of relying on FLOP counts alone.

\section*{Artifact Availability}

The artifact directory contains the original benchmark CSVs, the
\texttt{measured\_energy\_powermetrics} and
\texttt{measured\_energy\_powermetrics\_m4pro\_arch17\_cpu} directories with raw
\texttt{powermetrics} logs, per-window trial CSVs, summary CSVs, environment
metadata, measurement scripts, and the supplemental residualized GAN workflow.
The follow-up MPS-b1-224, CPU-b4-224, and CPU-b1-128-at-128px directories
provide the backend, batch-size, and input-size checks.
The retained M5 Pro summary should be paired with its raw artifact directory
before final camera-ready archival.

\begin{thebibliography}{00}

\bibitem{horowitz2014}
M. Horowitz, ``Computing's energy problem (and what we can do about it),''
in \emph{IEEE International Solid-State Circuits Conference Digest of
Technical Papers}, 2014.

\bibitem{eyeriss2016}
Y.-H. Chen et al., ``Eyeriss: An energy-efficient reconfigurable accelerator
for deep convolutional neural networks,'' \emph{IEEE Journal of Solid-State
Circuits}, vol. 52, no. 1, pp. 127--138, 2016.

\bibitem{mobilenetv3}
A. Howard et al., ``Searching for MobileNetV3,'' in \emph{Proc. IEEE/CVF
International Conference on Computer Vision}, 2019.

\bibitem{mobilenet}
A. G. Howard et al., ``MobileNets: Efficient convolutional neural networks
for mobile vision applications,'' arXiv:1704.04861, 2017.

\bibitem{mobilenetv2}
M. Sandler et al., ``MobileNetV2: Inverted residuals and linear
bottlenecks,'' in \emph{Proc. IEEE/CVF Conference on Computer Vision and
Pattern Recognition}, 2018.

\bibitem{efficientnet}
M. Tan and Q. Le, ``EfficientNet: Rethinking model scaling for convolutional
neural networks,'' in \emph{Proc. International Conference on Machine
Learning}, 2019.

\bibitem{squeezenet}
F. N. Iandola et al., ``SqueezeNet: AlexNet-level accuracy with 50x fewer
parameters and $<$0.5MB model size,'' arXiv:1602.07360, 2016.

\bibitem{shufflenet}
X. Zhang et al., ``ShuffleNet: An extremely efficient convolutional neural
network for mobile devices,'' in \emph{Proc. IEEE/CVF Conference on Computer
Vision and Pattern Recognition}, 2018.

\bibitem{greenai}
R. Schwartz et al., ``Green AI,'' \emph{Communications of the ACM}, vol. 63,
no. 12, pp. 54--63, 2020.

\bibitem{carbontracker}
L. F. W. Anthony, B. Kanding, and R. Selvan, ``Carbontracker: Tracking and
predicting the carbon footprint of training deep learning models,'' in
\emph{ICML Workshop on Challenges in Deploying and Monitoring Machine
Learning Systems}, 2020.

\bibitem{wgangp2017}
I. Gulrajani et al., ``Improved training of Wasserstein GANs,'' in
\emph{Advances in Neural Information Processing Systems}, 2017.

\end{thebibliography}

\end{document}
